\chapter{Apr.~28 --- GAGA}

\section{Revisiting Coherent Sheaves}

\begin{remark}
  Recall Chow's theorem says that if
  $X \subseteq \PP^n_\C$ is a closed analytic
  subvariety, then there exist
  $f_0, \dots, f_r \in \C[x_0, \dots, x_n]$
  homogeneous such that
  $X = V(f_0, \dots, f_r)$.
  Our goal is to compare
  coherent sheaves on
  $X^{\mathrm{alg}}$ and $X^{\mathrm{an}}$.
  This problem was solved by Serre
  in his GAGA paper.
\end{remark}

\begin{remark}
  Let $(X, \OO_X)$ be a ringed space,
  e.g. an algebraic variety or a complex analytic variety.
\end{remark}

\begin{definition}
  Let $\mathcal{F} \in \mathrm{Mod}_{\OO_X}$.
  \begin{enumerate}
    \item $\mathcal{F}$ is \emph{finite type}
      if at each $x \in X$, there exists
      $x \in U \subseteq X$ open and
      a surjection
      $\OO_U^{\oplus r} \twoheadrightarrow \mathcal{F}|_U$
      for some $r \ge 0$.
    \item $\mathcal{F}$ is
      \emph{coherent} if it is finite type
      and for every $U \subseteq X$ and
      morphism $\phi : \OO_U^{\oplus r} \to \mathcal{F}|_U$,
      we have that $\ker \phi$ is
      finite type (on $U$).
  \end{enumerate}
\end{definition}

\begin{remark}
  If $\mathcal{F} \in \mathrm{Mod}_{\OO_X}$ is coherent, then
  at each $x \in X$, there exists
  $x \in U \subseteq X$ open such that
  there exists an exact sequence
  \[
    \begin{tikzcd}
      \OO_U^{\oplus s}
      \ar[r] & \OO_U^{\oplus r}
      \ar[r] & \mathcal{F}|_U
      \ar[r] & 0
    \end{tikzcd}
  \]
\end{remark}

\begin{example}
  We have the following:
  \begin{enumerate}
    \item $\OO_X^{\oplus r}$ is finite type.
    \item If $(X, \OO_X)$ is an algebraic
      variety, then this new definition of
      coherence agrees with
      the old one.
  \end{enumerate}
\end{example}

\begin{theorem}[Oka-Cartan]\label{thm:oka-cartan}
  If $(X, \OO_X)$ is a complex analytic
  space, then $\OO_X$ is coherent.
\end{theorem}

\begin{proof}
  Look this up online, the proof uses
  the Weierstrass division theorem.
  Oka showed that $\OO_{\C^n}$ is coherent,
  and Cartan showed that if $X \subseteq \C^n$
  is a complex analytic subvariety, then
  $\mathcal{I}_X \subseteq \OO_{\C^n}$ is
  coherent (recall that $i_* \OO_X \cong \OO_{\C^n}/\mathcal{I}_X$ where $i : X \hookrightarrow \C^n$).
\end{proof}

\begin{remark}
  The category $\Coh_X$ of
  coherent sheaves on $X$ is an abelian
  category.
\end{remark}

\begin{definition}
  Let $\mathcal{F} \in \Coh_X$. Define
  \[
    H^i(X, \mathcal{F})
    = \varinjlim_{\mathcal{U}} H^i(\mathcal{U}, \mathcal{F}),
  \]
  where $\mathcal{U}$ ranges through all
  open (not necessarily affine) covers of $X$.
\end{definition}

\begin{remark}
  We have the following:
  \pagebreak
  \begin{enumerate}
    \item A short exact sequence in
      $\Coh_X$ gives rise to a longe
      exact sequences on $H^i$.
    \item If $(X, \OO_X)$ is a separated
      algebraic variety, then this new definition of
      $H^i$ agrees with the old one.
  \end{enumerate}
\end{remark}

\section{GAGA}

\begin{remark}
  Let $X \subseteq \PP^n_\C$ be a projective
  algebraic variety, and let
  $X^{\mathrm{an}} \subseteq \PP^{n, \mathrm{an}}_\C$ be
  $X$ viewed as a projective
  complex analytic variety.
  Note that there is a morphism of ringed
  spaces
  \[
    (X^{\mathrm{an}}, \OO_{X^{\mathrm{an}}})
    \longrightarrow (X, \OO_X)
  \]
  where $X^{\mathrm{an}} \to X$ is a set-theoretic
  bijection and
  $\OO_X \to i_* \OO_{X^{\mathrm{an}}}$ is
  given on open sets
  by $f \mapsto f^{\mathrm{an}}$.
  Letting $i^{-1}$ be the adjoint functor
  to $i_*$, we get a functor
  \begin{align*}
    \Phi : \Coh_X &\longrightarrow \Coh_{X^{\mathrm{an}}} \\
    \mathcal{F} &\longmapsto i^{-1} \mathcal{F} \otimes_{i^{-1} \OO_X} \OO_{X^{\mathrm{an}}},
  \end{align*}
  where the right-hand side is coherent
  by the Oka-Cartan theorem.
\end{remark}

\begin{remark}
  The functor $\Phi$ is exact. To prove
  this, Serre uses that the map
  \[
    \OO_{X, x}
    \longrightarrow \widehat{\OO}_{X, x}
  \]
  (here $\widehat{\OO}_{X, x}$ is the completion of $\OO_{X, x}$)
  is faithfully flat and
  $\widehat{\OO}_{X, x} \cong \widehat{\OO}_{X^{\mathrm{an}}, x}$.
\end{remark}

\begin{remark}
  If $\mathcal{U}$ is an open cover of $X$,
  then we get a map
  \[
    H^i(\mathcal{U}, \mathcal{F})
    \longrightarrow H^i(\mathcal{U}^{\mathrm{an}}, \Phi \mathcal{F}),.
  \]
  This induces a natural map (after taking
  direct limits)
  \[
    \phi^i_{\mathcal{F}} : H^i(X, \mathcal{F})
    \longrightarrow H^i(X^{\mathrm{an}}, \mathcal{F}^{\mathrm{an}}).
  \]
\end{remark}

\begin{theorem}[Theorem A]\label{thm:theorem-a}
  If $\mathcal{F} \in \Coh_X$, then
  $\phi^i_\mathcal{F}$ is an isomorphism.
\end{theorem}

\begin{proof}
  The sketch of the proof is as follows:
  \begin{enumerate}
    \item[0.] Reduce to the case $X = \PP^n$.
    \item Prove the result for
      $\mathcal{F} = \OO_{\PP^n}$.

      For this part, one notes that
      $H^0(\PP^n, \OO_{\PP^n}) = \C = H^0(\PP^{n, \mathrm{an}}, \OO_{\PP^{n, \mathrm{an}}})$ and
      for $i \ge 1$,
      \[
        H^i(\PP^n, \OO_{\PP^n})
        = 0 = H^i(\PP^{n, \mathrm{an}}, \OO_{\PP^{n, \mathrm{an}}}).
      \]
      The equality
      $H^i(\PP^{n, \mathrm{an}}, \OO_{\PP^{n, \mathrm{an}}}) = 0$ for $i \ge 1$
      follows from \emph{Hodge theory}.
    \item Prove the result for $\mathcal{F} = \OO_{\PP^n}(m)$.

      To prove this, use the short exact
      sequence (let $H = V(x_0)$)
      \[
        \begin{tikzcd}
          0 \ar[r] & \OO_{\PP^n}(m - 1)
          \ar[r, "\times x_0"]
               & \OO_{\PP^n}(m)
          \ar[r] & \OO_H(m) \ar[r] & 0
        \end{tikzcd}
      \]
      Then use induction
      on $n$ to see that
      $\phi^i_{\OO(m)}$ is an isomorphism
      if and only if $\phi^i_{\OO(m - 1)}$
      is.
    \item Prove the result for $\mathcal{F}$
      arbitrary.

      To do this, one chooses a presentation
      \[
        \begin{tikzcd}
          \bigoplus_{j = 1}^s \OO_{\PP^n}(n_j)
          \ar[r] & \bigoplus_{i = 1}^r \OO_{\PP^n}(m_i)
          \ar[r] & \mathcal{F} \ar[r] & 0
        \end{tikzcd}
      \]
      and uses the previous steps.
  \end{enumerate}
  See Serre's GAGA paper for more details.
\end{proof}

\begin{theorem}[Theorem B]\label{thm:theorem-b}
  For $\mathcal{F}, \mathcal{G} \in \Coh_X$,
  the map
  \[
    \Hom_{\OO_X}(\mathcal{F}, \mathcal{G})
    \longrightarrow \Hom_{\OO_{X^{\mathrm{an}}}}(\mathcal{F}^{\mathrm{an}}, \mathcal{G}^{\mathrm{an}})
  \]
  is an isomorphism.
\end{theorem}

\begin{proof}
  Set
  $\mathcal{A} := \Homm_{\OO_X}(\mathcal{F}, \mathcal{G})$ and
  $\mathcal{B} := \Homm_{\OO_{X^{\mathrm{an}}}}(\mathcal{F}^{\mathrm{an}}, \mathcal{G}^{\mathrm{an}})$.
  Then the result is equivalent to showing
  that $\Gamma(X, \mathcal{A})
    \to \Gamma(X^{\mathrm{an}}, \mathcal{B})$
  is an isomorphism. Consider the
  analytification of $\mathcal{A}$:
  \[
    \begin{tikzcd}
      \Gamma(X, \mathcal{A})
      \ar[r] \ar[dr, "\cong", swap] & {\Gamma(X^{\mathrm{an}}, \mathcal{B})} \\
                     & {\Gamma(X^{\mathrm{an}}, \mathcal{A}^{\mathrm{an}})} \ar[u]
    \end{tikzcd}
  \]
  The labeled map is an isomorphism by
  Theorem \ref{thm:theorem-a}, so it
  suffices to show $\mathcal{A}^{\mathrm{an}} = \mathcal{B}$.
  See GAGA.
\end{proof}

\begin{remark}
  Both Theorem \ref{thm:theorem-a} and Theorem \ref{thm:theorem-b}
  may be false if $X$ is not assumed to be
  projective. For example, consider
  $X = \Affine^1$ and
  $\mathcal{F} = \mathcal{G} = \OO_X$.
\end{remark}

\begin{theorem}\label{thm:theorem-c}
  For any $\mathcal{H} \in \Coh_{X^{\mathrm{an}}}$, there exists
  $\mathcal{F}$ such that
  $\mathcal{F}^{\mathrm{an}} \cong \mathcal{H}$.
  In particular,
  \[
    \Phi : \Coh_X \longrightarrow \Coh_{X^{\mathrm{an}}}
  \]
  is an equivalence of categories.
\end{theorem}

\begin{proof}
  The idea is to show
  $\mathcal{H}(m)$ is globally generated
  for $m \gg 0$. Use this to construct
  a presentation
  \[
    \begin{tikzcd}
      \bigoplus_{j = 1}^s \OO_{X^{\mathrm{an}}}(n_j)
      \ar[r, "\alpha"]
            & \bigoplus_{i = 1}^r \OO_{X^{\mathrm{an}}}(m_i)
      \ar[r] & \mathcal{H} \ar[r] & 0
    \end{tikzcd}
  \]
  and use Theorem \ref{thm:theorem-b} to
  show that $\alpha$ is algebraic.
\end{proof}

\begin{remark}
  Theorem \ref{thm:theorem-c} gives an
  alternative proof of Chow's theorem
  (Theorem \ref{thm:chow}). Fix an analytic
  subvariety $Z \subseteq \PP^n$. Then
  Cartan's theorem implies
  $\mathcal{I}_Z$ is coherent, and
  Theorem \ref{thm:theorem-c} implies
  there exists $\mathcal{F} \in \Coh_{\PP^n}$ such that
  $\mathcal{F}^{\mathrm{an}} \cong \OO_{\PP^{n, \mathrm{an}}}/\mathcal{I}_Z$.
  Now $\Supp \mathcal{F} = \Supp \mathcal{F}^{\mathrm{an}} = Z$.
\end{remark}
